% from ega4-8-fr.tex
\oldpage[IV-3]{51}
d'où une application canonique
\begin{equation}
\label{IV.8.13.4.2-fr}
\varprojlim_{\mu'}\left(\varinjlim_\mu
\operatorname{Hom}_S(X_\mu,X'_{\mu'})\right)
\longrightarrow
\operatorname{Hom}_S(\varprojlim_\mu X_\mu,
\varprojlim_{\mu'}X'_{\mu'}), \tag{8.13.4.2}
\end{equation}
évidemment fonctorielle en $\mathbf X$ et $\mathbf X'$, et qui complète la
définition du foncteur $L$.
\end{env}

\begin{proposition}[8.13.5]
\label{IV.8.13.5-fr}
Les hypothèses et notations étant celles de (8.13.4), le foncteur $L$ est
pleinement fidèle. Si de plus $S$ est un préschéma noethérien (ce qui implique
déjà que $S$ est quasi-compact et quasi-séparé (1.2.8)), $L$ est une
\emph{équivalence de catégories}.
\end{proposition}

Dire que $L$ est pleinement fidèle signifie que l'application (8.13.4.2) est
bijective quels que soient $\mathbf X$, $\mathbf X'$ dans $\mathcal C'$, ce
qui est un cas particulier de (8.13.2)~: en effet, les morphismes structuraux
$X_\mu\to S$ étant de présentation finie, sont en particulier quasi-compacts
et quasi-séparés, donc les $X_\mu$ sont quasi-compacts et quasi-séparés.

Pour montrer que lorsque $S$ est noethérien $L$ est une équivalence de
catégories, il suffit, puisque l'on sait déjà que $L$ est pleinement fidèle,
de prouver que tout $S$-préschéma essentiellement affine $X$ est
$S$-\emph{isomorphe} à un objet de la forme $L(\mathbf X)$ où
$\mathbf X\in\mathcal C'$ $(\mathbf 0_{\mathrm{III}},8.1.5)$.

Or, il y a par hypothèse une factorisation
$X\xrightarrow{g}X_0\xrightarrow{f}S$ du morphisme structural, $f$ étant de
présentation finie et $g$ affine. On peut donc écrire
$X=\operatorname{Spec}(\mathcal A)$, où $\mathcal A$ est une
$\mathcal O_{X_0}$-Algèbre quasi-cohérente (II, 1.3.1). Or, comme $X_0$ est
noethérien (puisqu'il en est ainsi de $S$ et que $f$ est de type fini),
$\mathcal A$ est limite inductive filtrante de la famille
$(\mathcal A_\lambda)$ de ses sous-$\mathcal O_{X_0}$-Algèbres quasi-cohérentes
de type fini (I, 9.6.6). Posons
$X_\lambda=\operatorname{Spec}(\mathcal A_\lambda)$~; les morphismes
$X_\lambda\to X_0$ sont de type fini, donc de présentation finie puisque
$X_0$ est noethérien, et par suite il en est de même des morphismes composés
$X_\lambda\to X_0\to S$ (1.6.2)~; autrement dit, les $X_\lambda$
appartiennent à $\mathcal C'_0$ et comme les morphismes $X_\lambda\to X_0$
sont affines, $\mathbf X=(X_\lambda)$ est un objet de $\mathcal C'$ dont la
limite projective existe et est $S$-isomorphe à $X$ en vertu de (8.2.2). Ceci
achève la démonstration.

\begin{remark}[8.13.6]
\label{IV.8.13.6-fr}
Il résulte de (1.6.2) et de (II, 1.6.2) que si $X$ et $Y$ sont
essentiellement affines au-dessus de $S$, alors il en est de même de
$X\times_S Y$. On en conclut par exemple $(\mathbf 0_{\mathrm{III}},8.2.5)$
qu'un $\mathcal C$-groupe n'est autre qu'un $(\mathbf{Sch})_{/S}$-groupe qui
est un préschéma essentiellement affine au-dessus de $S$. D'autre part, les
produits finis existent dans la catégorie $\mathcal C'$~: en effet, si
$\mathbf X=(X_\mu)_{\mu\in M}$, $\mathbf Y=(Y_\rho)_{\rho\in R}$ sont deux
objets de $\mathcal C'$, les produits $X_\mu\times_S Y_\rho$ sont des
$S$-préschémas de présentation finie, et en prenant pour morphismes de
transition
\[
X_\nu\times_S Y_\sigma\longrightarrow X_\mu\times_S Y_\rho
\]
les produits des morphismes de transition $X_\nu\to X_\mu$ et
$Y_\sigma\to Y_\rho$, on voit aussitôt que
$(X_\mu\times_S Y_\rho)$ est produit de $\mathbf X$ et de $\mathbf Y$ dans
$\mathbf{Pro}(\mathcal C'_0)$~; en outre (II, 1.6.2) les morphismes de
transition ainsi définis sont affines pour $\mu$ et $\rho$ assez grands, donc
le produit $\mathbf X\times\mathbf Y$ ainsi défini appartient bien à
$\mathcal C'$. On conclut alors comme plus haut qu'un $\mathcal C'$-groupe est
un $\mathbf{Pro}(\mathcal C'_0)$-groupe qui est essentiellement affine. On
déduit donc de (8.13.5) que les catégories des $\mathcal C'$-groupes et des
$\mathcal C$-groupes sont \emph{équivalentes} lorsque $S$ est
\emph{noethérien}. Il semble plausible que lorsque $S$ est le spectre d'un
corps $k$, la catégorie des $\mathcal C$-groupes soit \emph{équivalente} à
celle des $\boldsymbol K$-groupes, où $\boldsymbol K$ est la catégorie des
$S$-préschémas \emph{quasi-compacts}~; autrement dit, tout préschéma en
groupes sur $k$ qui est quasi-compact serait essentiellement affine. D'autre
part, si l'on désigne
